\documentclass{smartmeter-report}   % add [final] to fail the build on any \placeholder

\projecttitle{Power Consumption Forecasting Using Smart Meter Data}
\shorttitle{Smart Meter Forecasting}
\phase{1}
\phasename{Problem Definition \& Data Source Identification}
\deliverable{2--3 page project proposal}
\submissiondate{September 29, 2026}
\team{\placeholder{Team member 1}, \placeholder{Team member 2}, \placeholder{Team member 3}}
\course{Data Analytics, UNB, Fall 2026}

\begin{document}
\maketitle

% Brief: 2-3 pages excluding references/appendices. Headings must match the five
% proposal components below. Cite external data sources. Clearly separate CONFIRMED
% data access from PLANNED or potential sources.

\section{Problem Statement}
% Real-world problem, who is affected, why it matters, the system studied,
% stated specifically and measurably (not a broad topic).
\placeholder{Households have little visibility of what their electricity use and bill
will be in the coming weeks. Describe the problem of forecasting household consumption
(and the resulting bill) from half-hourly smart meter readings for London households,
who benefits (households, suppliers, network operators), and state it measurably, e.g.\
forecast next-period consumption per household within a stated error.}

\section{Project Objectives \& Analytics Goal}
% 2-4 specific objectives; primary analytics task; target variable; what success enables.
\begin{enumerate}
  \item \placeholder{Objective 1, e.g.\ build a reproducible pipeline from raw readings to a clean household panel}
  \item \placeholder{Objective 2, e.g.\ forecast each household's consumption over a fixed horizon}
  \item \placeholder{Objective 3, e.g.\ convert forecast consumption into a projected bill}
\end{enumerate}

\textbf{Primary analytics task:} Prediction (regression on future consumption).\\
\textbf{Target variable:} \placeholder{forecast horizon and granularity, e.g.\ daily
consumption in \kwh{} per household over the next 7 days}; projected bill derived from it.\\
\textbf{Success looks like:} \placeholder{what the project will produce or enable, with the
baseline it must beat}.

\section{Data Source Identification}
% For each source: contents, relevance, type, format, size / period.
The project uses the \enquote{Smart meters in London} dataset \parencite{kaggle_smartmeters},
a refactored release of the Low Carbon London trial data published by UK Power Networks
\parencite{ukpn_lcl}, together with the weather and household metadata distributed with it.

\begin{table}[H]
\centering\small
\smrows
\begin{tabularx}{\linewidth}{L{3.6cm} L{2.2cm} Y L{2.6cm}}
\headrow \hd{Source} & \hd{Type} & \hd{What it provides} & \hd{Access / feasibility} \\
Half-hourly readings (\code{halfhourly_dataset/}, 112 CSVs)
  & Sensor; public/open
  & Energy use in \kwhhh{} for \num{5566} households, Nov 2011 -- Feb 2014; \num{167817021} rows
  & Downloaded; confirmed \\
Daily aggregates (\code{daily_dataset/})
  & Derived; public/open
  & Per-household daily mean, median, max, min, std and total energy
  & Downloaded; confirmed \\
Household information (\code{informations_households.csv})
  & Public/open
  & Tariff (standard vs.\ time-of-use) and Acorn socio-demographic group per household
  & Downloaded; confirmed \\
Acorn profiles (\code{acorn_details.csv})
  & External
  & Demographic and lifestyle indices for each Acorn category \parencite{caci_acorn}
  & Downloaded; confirmed \\
Weather (\code{weather_hourly_darksky.csv}, \code{weather_daily_darksky.csv})
  & External; API
  & Hourly and daily London temperature, humidity, wind, pressure \parencite{darksky}
  & Downloaded; confirmed \\
UK bank holidays (\code{uk_bank_holidays.csv})
  & Public/open
  & Holiday calendar for the study period
  & Downloaded; confirmed \\
\placeholder{Tariff price schedule}
  & External
  & Unit prices needed to turn forecast \kwh{} into a bill amount
  & \placeholder{Planned; source to confirm} \\
\end{tabularx}
\caption{Data source checklist.}
\label{tab:sources}
\end{table}

\section{Data Access, Quality \& Feasibility}
% How the data is obtained; known issues; privacy/licensing/ethics; sufficiency.
\subsection{Access}
\placeholder{All sources above except the tariff prices are downloaded and profiled.
State how they were obtained and their licence terms.}

\subsection{Known data-quality issues}
% Figures below come from checks/profile_report.md; re-run the profiler if the data changes.
An initial profile of the half-hourly readings found:
\begin{itemize}
  \item \num{5560} readings recorded as the string \code{Null}, all at off-grid timestamps
        (mostly on 18 December 2012).
  \item \SI{0.27}{\percent} of expected half-hour slots are missing; 95\,\% of households
        have at least one gap, the longest spanning about 210 days.
  \item Staggered coverage: households join between November 2011 and October 2013;
        179 households have less than a year of data.
  \item 231 households contain at least one full day of consecutive zero readings,
        suggesting meter faults or vacant homes.
\end{itemize}

\subsection{Privacy and ethics}
\placeholder{Households are pseudonymised (\code{LCLid}); no personal data; note any
licence or ethical constraints.}

\subsection{Sufficiency}
\placeholder{Why these sources are enough, or what extra data might be needed.}

\section{Preliminary Project Scope}
% In / out of scope; stakeholders; expected final outputs.
\textbf{In scope:} \placeholder{\ldots}\\
\textbf{Out of scope:} \placeholder{\ldots}\\
\textbf{Stakeholders:} \placeholder{\ldots}\\
\textbf{Expected outputs:} \placeholder{e.g.\ forecasting model, analytical report, dashboard}

\printbibliography[heading=bibintoc]

\end{document}
